We present a systematic study on laboratory mouse detection using the
PaddleDetection~2.6 framework on a dual Tesla T4 GPU server.
Starting from three heterogeneous data sources, we construct a
7,286-image VOC-format dataset (5,828 train / 1,458 val) covering
two classes: \textit{mouse} and \textit{other}.
We design a $4\!\times\!2$ controlled-variable experiment matrix
spanning two model families (PicoDet-S and YOLOv3-MobileNetV1),
two GPU configurations (single / dual card), and two data scales
(full dataset / one-third subset), yielding eight training runs
totalling approximately 270 GPU-hours.

Learning rates are derived from the Linear Scaling Rule
($\eta = \eta_\text{base}\!\times\! B/B_\text{ref}$) to isolate
the effect of batch size from parallelism.
Our best model, PicoDet-S with single-GPU small-batch training
(batch\,=\,32, lr\,=\,0.08), achieves \textbf{mAP@0.5 = 94.30\%}
at epoch~70, outperforming YOLOv3's best result
(91.54\%) by \textbf{+2.76 percentage points} while being
\textbf{21$\times$ smaller} (4.4\,MB vs.\ 92\,MB) and running at
\textbf{2.2$\times$ higher FPS} on the same GPU.
A counter-intuitive finding is that the small-batch single-GPU run
outperforms the standard dual-GPU run (+0.5\,pp), which we attribute
to gradient noise acting as implicit regularisation on the
5,828-sample dataset.

The trained PicoDet-S model is exported to ONNX, integrated into a
React Native application with ONNX Runtime and CoreML back-end,
and deployed on an iPhone.  Through four engineering iterations,
inference speed improves from 1\,FPS to a stable \textbf{14\,FPS}
in warm-machine conditions.
We also study the dramatic effect of dataset scale on YOLOv3:
tripling the training set (1,942~$\to$~5,828 images) lifts mAP
by \textbf{+47 percentage points} (44.1\%~$\to$~91.5\%),
confirming that data engineering is the primary driver of accuracy
for anchor-based heavy architectures on small-domain datasets.
